% BHU PYQ -- Manually transcribed
% Paper: MTB-302 : Differential Equations (New & Old Course)
% Exam: B.Sc. (Hons.) Semester III Examination, 2016-17

\documentclass[11pt,a4paper]{article}
\usepackage[a4paper,top=2.5cm,bottom=2.5cm,left=2.8cm,right=2.8cm,headheight=14pt]{geometry}
\usepackage{amsmath,amssymb}
\usepackage[T1]{fontenc}
\usepackage[utf8]{inputenc}
\usepackage{lmodern,microtype,fancyhdr,enumitem,hyperref,lastpage,parskip}

\hypersetup{colorlinks=true,urlcolor=black,linkcolor=black,
  pdftitle={MTB-302 Differential Equations -- BHU B.Sc. Sem III 2016-17}}

\pagestyle{fancy}\fancyhf{}
\renewcommand{\headrulewidth}{0.4pt}\renewcommand{\footrulewidth}{0.4pt}
\fancyhead[L]{\small   \textit{Banaras Hindu University}}
\fancyhead[R]{\small   \textit{MTB-302: Differential Equations}}
\fancyfoot[C]{\small Page  \thepage\ of \pageref{LastPage}}


\newlist{parts}{enumerate}{2}
\setlist[parts,1]{label=(\alph*),leftmargin=*,itemsep=5pt,topsep=3pt}
\setlist[parts,2]{label=(\roman*),leftmargin=*,itemsep=3pt,topsep=2pt}

\newcommand{\pts}[1]{\hfill   \textit{[#1]}}


\begin{document}

\begin{center}
  {\large\bfseries BANARAS HINDU UNIVERSITY}\\[2pt]
  {\normalsize B.Sc.\ (Hons.) --- Semester III Examination, 2016-17}\\[6pt]
  \rule{0.8\linewidth}{0.5pt}\\[6pt]
  {\large\bfseries Paper No.\ MTB-302: Differential Equations  extnormal{(New \& Old Course)}}\\[4pt]
  {\normalsize Time: 3 Hours \hfill Maximum Marks: 70}
\end{center}
\rule{\linewidth}{0.4pt}
\smallskip



\noindent   \textit{Instructions: Answer   \textbf{five} questions including Question~1 which is compulsory. Figures in right-hand margin indicate marks.}
\bigskip




\noindent   \textbf{Question 1.}   \textit{(Compulsory --- 2 marks each.)}

\begin{parts}
  \item If $y_1, y_2, y_3$ are solutions of $y''+ay'+by=0$, show that the Wronskian $W(y_1,y_2,y_3)=0$.
  \item Show that the number of integrating factors of $M(x,y)\,dx+N(x,y)\,dy=0$, which has a solution, is infinite.
  \item Find the general solution of $p^2+p-6=0$, where $p=\dfrac{dy}{dx}$.
  \item Define singular solution of a differential equation.
  \item Obtain the linear homogeneous 2nd order ODE whose linearly independent solutions are $y=x$ and $y=x^2$.
  \item Transform $x^2\dfrac{d^2y}{dx^2}+ax\dfrac{dy}{dx}+by=0$ by substituting $x=e^t$.
  \item Obtain the condition for $y=e^{mx}$ to be a solution of $\dfrac{d^2y}{dx^2}+c\dfrac{dy}{dx}+Py=0$.
  \item Examine whether $\sin x\,\dfrac{dy}{dx}+(\cos x+\sin x)\dfrac{dy}{dx}+y\cos x=0$ is exact.
  \item Define regular singular points for a 2nd-order linear ODE and give an example.
  \item Write the value of $P_n(0)$ for $n=0,1,2,\ldots$
  \item Show that $J_{1/2}(x)=\sqrt{\dfrac{2}{\pi x}}\sin x$.
\end{parts}
\medskip\rule{\linewidth}{0.2pt}

\medskip


\noindent   \textbf{Question 2.}

\begin{parts}
  \item Show that the ODE $M\,dx+N\,dy=0$ (where $M,N$ have continuous partial derivatives) is exact if and only if $\dfrac{\partial M}{\partial y}=\dfrac{\partial N}{\partial x}$.\pts{6}
  \item Solve: $(x+y+1)\,dx=(2x+2y+1)\,dy$.\pts{6}
\end{parts}
\medskip\rule{\linewidth}{0.2pt}

\medskip


\noindent   \textbf{Question 3.}

\begin{parts}
  \item Solve $\dfrac{d^2y}{dx^2}+y=\cos x$ by the method of variation of parameters.\pts{6}
  \item Solve $x^2\dfrac{d^2y}{dx^2}-2x(3x-2)\dfrac{dy}{dx}+3x(3x-4)y=e^{3x}$ by reducing to normal form.\pts{6}
\end{parts}
\medskip\rule{\linewidth}{0.2pt}

\medskip


\noindent   \textbf{Question 4.}

\begin{parts}
  \item Find the general solution of $n$th-order linear homogeneous ODE with constant coefficients when the auxiliary equation has real and distinct roots.\pts{6}
  \item Find the solution of $\dfrac{d^3y}{dx^3}-5\dfrac{d^2y}{dx^2}+\dfrac{dy}{dx}+y=xe^x+e^x$ satisfying $y=0$, $\dfrac{dy}{dx}=0$, $\dfrac{d^2y}{dx^2}=0$ at $x=0$.\pts{6}
\end{parts}
\medskip\rule{\linewidth}{0.2pt}

\medskip


\noindent   \textbf{Question 5.}

\begin{parts}
  \item Reduce $(px^2+y^2)(px+y)=(p+1)^2$ (where $p=dy/dx$) to Clairaut's form using $u=xy$, $v=x+y$ and find the complete primitive.\pts{6}
  \item Find the singular solution of $\sin\!\left(x\dfrac{dy}{dx}\right)\cos y - \cos\!\left(x\dfrac{dy}{dx}\right)\sin y = 0$.\pts{6}
\end{parts}
\medskip\rule{\linewidth}{0.2pt}

\medskip


\noindent   \textbf{Question 6.}

\begin{parts}
  \item Solve $\dfrac{d^2y}{dx^2}+(x^2-1)\dfrac{dy}{dx}+y=(e^x-1)y$ in series about the ordinary point $x=0$.\pts{6}
  \item Express $f(x)=x^4+3x^3-x^2+5x-2$ in terms of Legendre polynomials.\pts{6}
\end{parts}
\medskip\rule{\linewidth}{0.2pt}

\medskip


\noindent   \textbf{Question 7.}

\begin{parts}
  \item Show that $\dfrac{d}{dx}\left[x^n J_n(x)\right]=x^n J_{n-1}(x)$.\pts{6}
  \item Prove that $P_n(x)=\dfrac{1}{2^n n!}\dfrac{d^n}{dx^n}(x^2-1)^n$, where $P_n(x)$ is the Legendre polynomial of degree $n$.\pts{6}
\end{parts}

\vfill\rule{\linewidth}{0.4pt}

\begin{center}\small
\href{https://bhu-exam.vercel.app/}{Click here to visit our website}\\[4pt]
\href{https://me-aryan.vercel.app/}{Click here to visit our contributor}\\[4pt]
\href{https://quashan-me.vercel.app/}{Click here to visit our contributor}
\end{center}
\end{document}
